% Glossary terms from ch03: lines of the form \item[Term] Definition.
\item[Dijkstra's algorithm] The algorithm that computes the cheapest cost from one source to every vertex of a graph with non-negative edge costs by repeatedly settling the open vertex with the smallest tentative cost.
\item[Settled vertex] A vertex that has been popped from the priority queue and added to the closed set; its tentative cost is then final.
\item[Relaxation] The update of a neighbour's tentative cost and parent pointer when a cheaper path to it through the vertex being expanded is found.
\item[Tentative cost] The cost $\gcost(v)$ of the best path to $v$ found so far; it never underestimates the true distance and only ever decreases.
\item[Shortest-path distance] The smallest cost $\delta(u,v)$ of any path from $u$ to $v$, taken to be infinite when no path exists.
\item[Shortest-path tree] The tree formed by the parent pointers of a single-source search; following them from any vertex back to the source traces a shortest path.
\item[Expansion] One non-stale pop from the priority queue together with the generation of the successors of the popped vertex; the unit in which the work of a search is counted.
\item[Stale entry] A priority-queue entry of a vertex that has already been settled with a smaller key; lazy deletion discards it when it is popped.
\item[Uniform-cost search] Dijkstra's algorithm with an early exit: the search stops as soon as the target is popped, and only vertices no farther than the target are settled.
\item[Early exit] Stopping a search when the target is popped from the queue, never when it is first discovered.
\item[Backward Dijkstra] A run of Dijkstra's algorithm from the goal over the graph with every edge reversed, which yields the distance to the goal from every vertex.
\item[True-distance heuristic] The exact cost-to-go $\hcost^*(v)=\delta(v,t)$ computed by backward Dijkstra; it is consistent, exact on the static map and still admissible in space-time and under constraints.
\item[Multi-source Dijkstra] Dijkstra's algorithm started with several sources at cost zero, which yields the cost to the nearest source for every vertex.
\item[Bidirectional search] Two simultaneous waves, one forward from the source and one backward from the target, stopped when the sum of the two smallest queue keys reaches the cost of the best connecting path found.
\item[Breadth-first search] The unit-cost special case of Dijkstra's algorithm in which a first-in-first-out queue replaces the heap.
